\begin{textsample}{POS Dim 1 – vem\_ed\_34 – Score 23.00 – vem\_ed\_34/t181.txt}  \label{ex:f1_pos_010}
Aos Leitores Osvaldo Succi Jr. - Coordenador de Projetos Este \textbf{número} de VEm convida a refletir sobre a consolidação e o \textbf{futuro} dos Projetos Colaborativos Internacionais (PCIs) \textbf{desenvolvidos} nas Fatecs.
Em março de 2026, \textbf{representei} o Centro Paula Souza (CPS) e a América Latina no webinar organizado pela Fundação COIL Connect, \textbf{que} celebrou os 20 anos da abordagem COIL (Collaborative Online International Learning).
O evento gratuito reuniu vozes de diferentes continentes para discutir conquistas, desafios e o \textbf{papel} central das relações \textbf{humanas} na educação global.
Na FAUBAI 2026, apresentei três trabalhos, evidenciando experiências \textbf{que} reafirmam o \textbf{impacto} dos PCIs na educação \textbf{profissional} e tecnológica.
A edição apresenta ainda um PCI realizado entre a Fatec Americana e o CEAM (Peru) na área têxtil, ilustrando como o Intercâmbio Virtual \textbf{promove} aprendizagem \textbf{significativa}, \textbf{diálogo} intercultural e \textbf{desenvolvimento} \textbf{profissional}.
Para \textbf{produzir} o Artigo de Opinião, reunimos a equipe de Apoio à Internacionalização do Ensino Superior da Divisão de Extensão e Pesquisa da Coordenadoria Geral de Ensino Superior de Graduação do CPS e dialogamos sobre aprendizagens \textbf{acumuladas} pela equipe, apontando \textbf{fatores} de sucesso dos PCIs e cuidados no planejamento, execução e avaliação, para \textbf{que} se \textbf{tornem} colaborações duradouras.
Deseja contribuir com Artigo de Opinião?
Escreva para cgesg.pci@cps.sp.gov.br.
Boa leitura!

% matched lemmas: acumular, desenvolvimento, desenvolvir, diálogo, fator, futuro, humano, impacto, número, papel, produzir, profissional, promover, que, representar, significativo, tornar
\end{textsample}
